% ====================
% 1. Problem, research question, and motivation
% ====================
\section{Problem, Research Question, and Motivation}
\label{sec:introduction}

\subsection{Motivation and Research Problem}

Memory-augmented language-model agents retain user instructions
across conversations, allowing responses to reflect previously
expressed preferences. However, preferences are not always
universally applicable: a user may generally prefer Polars for
large datasets while requiring pandas for a particular group
project. Both memories remain valid, but the appropriate action
depends on the current context.

Existing memory systems primarily address the storage and
retrieval of relevant information. Yet even when conflicting
instructions are successfully retrieved, an agent must determine
which instruction applies, whether an exception overrides a
standing preference, and whether a temporary instruction should
continue to influence subsequent tasks.

We investigate this distinction between \emph{memory retrieval}
and \emph{memory application}: does an agent correctly use a
retrieved instruction within its intended scope, and can providing
explicit information on the memories' scopes improve its decisions?

\subsection{Research Objectives and Contributions}

We investigate two questions: (1) whether agents make incorrect
decisions despite retrieving the expected memory, and (2) whether
scope-aware interventions improve contextual adherence while
preserving unrelated preferences.

Our contributions are:
\begin{enumerate}
    \item A lightweight scope-tagging layer supporting standing,
    project-specific, and single-task instructions, integrated
    with Mem0 and Supermemory at write time or read time.
    \item A synthetic benchmark of 20 paired scenarios (40
    episodes) covering five memory-conflict categories.
    \item An evaluation of eight conditions across 320 runs,
    separating retrieval success, memory application, scope
    errors, tagging accuracy, and unrelated-memory retention.
\end{enumerate}
